\documentclass[10pt,a4paper]{amsart}
\usepackage[margin=19mm]{geometry}
\usepackage{amsmath,amssymb,mathtools}
\usepackage[scr=rsfs,cal=euler]{mathalfa}
\usepackage{xfrac}
\usepackage[unicode,hidelinks,bookmarks=false]{hyperref}
\newcommand{\ie}{i.e.\ }
\newcommand{\cf}{cf.\ }
\newcommand{\resp}{respectively\ }
\newcommand{\Subsection}{\subsection}
\providecommand{\nobreakdash}{\nobreak\hbox{-}}
\setlength{\emergencystretch}{2em}
\multlinegap0pt
\usepackage{amssymb}
\usepackage[shortlabels]{enumitem}
\usepackage[normalem]{ulem}
\newtheorem{lemma}{Lemma}
\newtheorem{proposition}{Proposition}
\usepackage{cleveref}
\crefname{lemma}{Lemma}{Lemmas}
\crefname{proposition}{Proposition}{Propositions}
\newcommand{\Fill}{\operatorname{Fill}}
\newcommand{\mc}{\mathcal}
\newcommand{\mr}{\mathring}
\newcommand{\psub}{\mr W}
\title{On fixed points of pseudo-Anosov maps}
\author{Tarik Aougab, David Futer, Samuel J. Taylor}
\date{Source: arXiv:2509.07818v2 (CC BY 4.0)}
\begin{document}
\maketitle

\noindent\emph{Source and licence.} On fixed points of pseudo-Anosov maps. Version arXiv:2509.07818v2 (CC BY 4.0), distributed by arXiv under the Creative Commons Attribution 4.0 International licence, http://creativecommons.org/licenses/by/4.0/, as recorded in the arXiv licence metadata for this exact version and shown on the abstract page https://arxiv.org/abs/2509.07818v2. Full licence notice: \texttt{licenses/arxiv-2509.07818-CC-BY-4.0.txt}.\par\smallskip

\noindent\textit{Editorial scope.} Counted unit: the article's proposition Prop:FromTriangToEdge (source0.tex lines 1817--1825). The article's complete original proof of it is delivered with it (source0.tex lines 1828--1867, one \texttt{proof} environment of 40 lines). No mathematics is written, repaired or re-derived: the statement and the proof are the article's own lines, in the article's own notation, and no retained line is substituted or re-flowed.  Retained uncounted as the source-local context the proof needs: lines 1037--1040; lines 1313--1324; lines 1725--1731; lines 1792--1798.  Nothing was omitted from any retained span, so the package carries no omission record.  No retained line contains a citation, so the wrapper carries no bibliography and no bibliography entry.  No retained line contains a figure environment, an image-inclusion command or drawing code, so no image is delivered and the package declares no asset dependency.  Editorial formatting only, in addition: an AMS \texttt{amsart} preamble that restates the article's own declarations and macros used by the retained lines, namely the environments \texttt{lemma}, \texttt{proposition} on separate counters, together with the standard packages those lines need.  The retained environments print in this document as Lemma 1, Proposition 1 in order of appearance, whereas the article prints the same items with the numbering of its own full document.  The complete native source of the article as distributed in the arXiv e-print for this exact version is bundled unchanged as \texttt{sources/184777/source0.tex}.
\begin{lemma}[Cutting subsurfaces]\label{Lem:CuttingSubsurfaces}
Let $\psub$ be a proper essential subsurface of $\mr S$ with $d_{\psub} (\lambda^+, \lambda^-) \geq 4$. 
Let $Y = \Fill_S(\psub)$, and suppose that $f$ overlaps $Y$. Then, for any section $T$, there is an edge $\sigma$ such that both $\sigma$ and $f(\sigma)$ cut $\psub$. 
\end{lemma}

\begin{lemma}
%[Closest sections; subsurface distance]
\label{Lem:ClosestSecSubsurfaceDist}
Let $T$ be a section of the veering triangulation and $\sigma$ an edge of $T$. 
Assume that $i(\sigma, f(\sigma)) > 0$. Let $P = P(\sigma, f(\sigma))$ be the associated pocket, with top and bottom sides $P^+$ and $P^-$.
Then
for any essential subsurface $\psub \subset \mr S$,
\[
d_{\psub}(\sigma , f(\sigma)) \: \leq \: d_{\psub}(P^+, P^-) \: \leq \:  d_{\psub}(\sigma , f(\sigma)) + 9.
\]
% where an undefined projection is interpreted to have $0$ distance.
\end{lemma}

\begin{proposition}[Choi--Rafi]\label{Prop:ChoiRafi}
For any pair of markings $\mu_1,\mu_2$ on $S$,
\[
\log i(\mu_1,\mu_2) \asymp d_S (\mu_1, \mu_2) + \sum_Y \big[ d_Y(\mu_1,\mu_2) \big] + \sum_A \big[ \log \left (d_A(\mu_1,\mu_2) \right ) \big],
\]
where the first sum is over proper non-annular subsurfaces of $S$, and the second sum is over annuli. 
\end{proposition}

\begin{lemma}\label{Lem:SSumSigma}
Let $f$ be a pseudo-Anosov, and let $\sigma$ be any edge of the veering triangulation. Then
\[
\log i(\sigma, f(\sigma)) \asymp \mc S(\sigma, f(\sigma)),
\]
with the convention that $\log(0)=0$.
\end{lemma}

\begin{proposition}\label{Prop:FromTriangToEdge}
Let $f$ be strongly irreducible, and let $T$ be a section. 
Let $\sigma$ be an edge of $T$ for which the intersection number $i(\sigma, f(\sigma))$ is largest. Then
% Then there is an edge $\sigma \subset T$ such that
\[
\log i(\sigma, f(\sigma)) \asymp \log i(T, f(T)),
\]
with the convention that $\log(0)=0$.
\end{proposition}

\begin{proof}
%    Fix an $f$--section $T$, and let $\sigma$ be the edge of $T$ for which the sum $\mc S(\sigma, f(\sigma))$ is largest. 
Let $\sigma'$ be an arbitrary edge of $T$. By \Cref{Lem:SSumSigma}, combined with the hypothesis that $\sigma$ is the edge of $T$ that maximizes the intersection number $i(\sigma, f(\sigma))$, we have
\[
\mc S(\sigma', f(\sigma'))  \asymp  \log i(\sigma', f(\sigma')) \leq 
 \log i(\sigma, f(\sigma)) \asymp  \mc S(\sigma, f(\sigma)).
\]
Comparing the first and last terms, we obtain
\begin{equation}\label{Eqn:SigmaLargest}
\mc S(\sigma', f(\sigma')) \prec \mc S(\sigma, f(\sigma))
\end{equation}
for every $\sigma' \subset T$. 


Next, we relate $\mc S (\sigma, f(\sigma))$ to $\mc S(T, f(T))$ by a partitioning argument, as follows. 
Since $f$ is strongly irreducible, every subsurface $\mr Y$ is overlapped by $f$.
Thus, by  \Cref{Lem:CuttingSubsurfaces}, 
% and the hypothesis that $f$ is strongly irreducible, 
every subsurface $\mr Y$ that contributes to the sum $\mc S(T, f(T))$ must be cut by both $\sigma'$ and $f(\sigma')$ for \emph{some} edge $\sigma'$ of $T$, depending on $\mr Y$. Hence, we have $d_{\mr Y}(T, f(T)) = d_{\mr Y}(\sigma', f(\sigma')) + O(1)$ for 
this choice of $\sigma'$.
Then we obtain
%    , applying \Cref{Lem:ClosestSecSubsurfaceDist} to each summand gives
\[
\mc S(\sigma, f(\sigma)) \leq \mc S(T, f(T)) \prec \sum_{\sigma' \subset T} \mc S(\sigma', f(\sigma'))
\prec \mc S(\sigma, f(\sigma)).
\]
Here, the first inequality is just the fact that $\sigma \subset T$. The next (coarse) inequality comes from decomposing $\mc S(T, f(T))$ into a sum over edges: for each edge $\sigma'$, let $\mc S(\sigma', f(\sigma'))$ be the sub-sum of $\mc S(T, f(T))$ such that $\sigma'$ and $f(\sigma')$ cut the subsurface in question. Since each subsurface with an associated summand in $\mc S(T, f(T))$ contributes to a summand in at least one such $\mc S(\sigma', f(\sigma'))$, the second inequality follows. The final coarse inequality uses the fact that there are at most $6|\chi(S)|$ edges of $T$, 
combined with \eqref{Eqn:SigmaLargest}.
% the hypothesis that $\sigma$ is the edge of $T$ for which the sum $\mc S(\sigma, f(\sigma))$ is largest. 

Since the first and last terms of the above equation are equal, we obtain
\[
\mc S(\sigma, f(\sigma)) \asymp \mc S(T, f(T)).
\]
Putting it all together gives
\[
\log i(\sigma, f(\sigma)) \asymp \mc S(\sigma, f(\sigma)) \asymp \mc S(T, f(T)) \asymp \log i(T, f(T)), 
\]
where the last coarse equality is again \Cref{Prop:ChoiRafi}.
\end{proof}

\end{document}
